\documentclass[a4paper,twoside,titlepage]{article}

% Informazioni sull'autore e sul corso
\newcommand{\Autore}{Edoardo~Porcaro}
\newcommand{\MailAutore}{\href{mailto:ufogamespa@gmail.com}{\Autore{}}}
\newcommand{\NomeCorso}{<!-- TITOLO CORSO QUI -->}
\newcommand{\NomeProf}{<!-- NOME PROF QUI -->}
\newcommand{\AnnoAccademico}{2025/2026}

% Carica pacchetto personale
\input{porcaro}

% Titolo, autore e data
\title{\sffamily \NomeCorso\\\vspace{6pt}\Large Prof.~\NomeProf\vspace{-8pt}}
\author{\sffamily Appunti di \Autore}
\date{\sffamily Anno accademico \AnnoAccademico}

% -------[ DEFINIZIONI PERSONALIZZATE ]--------------------------------------------- 
    %  <!-- QUI -->
% -------[ fine ]-------------------------------------------------------------------

% -------[ ACRONIMI ]---------------------------------------------------------------
    % <!-- QUI -->
    % \newacronym{api}{API}{Application Program Interface} --> \gls{api}
% -------[ fine ]-------------------------------------------------------------------

% -------[ INIZIO DOCUMENTO ]-------------------------------------------------------
\begin{document}

% Copertina
\maketitle
\thispagestyle{empty}
\cleardoublepage

% Introduzione
\begin{abstract}
    % <!-- TESTO DEL RIASSUNTO -->
\end{abstract}
\thispagestyle{empty}
\cleardoublepage

% Indice dei contenuti e delle figure
\tableofcontents
\clearpage
%\listoffigures
%\clearpage
%\listoftables
%\clearpage

% Contenuto del documento
\input{contenuto}

% Appendici
\begin{appendices}
    % Appendice acronimi
    \section[Acronimi]{Acronimi}
    \printglossary[type=\acronymtype,title={Acronimi}]
\end{appendices}

\end{document}
% -------[ FINE DOCUMENTO ]---------------------------------------------------------
